\documentclass[../Main.tex]{subfiles}

\begin{document}


\section{Verzeichnisse}

\subsection{Literatur- und Quellenverzeichnis}

\printbibliography[heading=none]

\subsection{Abbildungsverzeichnis}

\listoffigures

\subsection{KI-Deklaration}

Im Rahmen dieser Arbeit wurden KI-gestützte Werkzeuge unterstützend eingesetzt. Ihr Einsatz beschränkte sich auf die Fehlersuche im Quellcode (Debugging), die Identifikation möglicher Probleme in frühen Entwicklungsphasen sowie die sprachliche Überprüfung des Begleitberichts.

Wurde bei der Entwicklung einer Quelldatei KI oder ein anderes externes Hilfsmittel verwendet, so ist dies im Header der entsprechenden Datei dokumentiert. Sämtliche zur Bewertung eingereichten Codeabschnitte wurden von den Autoren selbst entwickelt. Auch alle Abbildungen und Grafiken wurden ohne den Einsatz generativer KI erstellt.

KI-Systeme wurden nicht zur Erstellung inhaltlicher Bestandteile dieser Arbeit verwendet. Insbesondere stammen die Forschungsfrage, die Analysen, die technischen Konzepte, die Schlussfolgerungen sowie die schriftlichen Inhalte des Begleitberichts von den Autoren. Der Einsatz von KI beschränkte sich auf sprachliche Korrekturen, formale Qualitätskontrollen sowie die Überprüfung, ob die im Projektvertrag definierten Anforderungen eingehalten wurden.

\end{document}
